\documentclass[11pt]{article}
\usepackage[margin=0.9in]{geometry}
\usepackage{amsmath,amssymb,amsthm}
\usepackage{booktabs,enumitem,multicol,xcolor,fancyhdr}
\usepackage[most]{tcolorbox}
\usepackage{listings}
\usepackage{pdflscape}

\definecolor{accent}{HTML}{1F4E79}
\definecolor{codebg}{HTML}{F5F5F5}
\definecolor{tipbg}{HTML}{FFF6DD}

\lstset{
  language=Python, basicstyle=\ttfamily\small, backgroundcolor=\color{codebg},
  keywordstyle=\color{accent}\bfseries, commentstyle=\color{gray}\itshape,
  frame=single, rulecolor=\color{gray!40}, columns=fullflexible,
  keepspaces=true, showstringspaces=false, aboveskip=6pt, belowskip=6pt,
  morekeywords={np}
}

\newtcolorbox{keybox}[1][]{colback=accent!6, colframe=accent, boxrule=0.6pt,
  arc=2pt, left=6pt, right=6pt, top=4pt, bottom=4pt, fonttitle=\bfseries, #1}
\newtcolorbox{tipbox}[1][]{colback=tipbg, colframe=orange!70!black, boxrule=0.6pt,
  arc=2pt, left=6pt, right=6pt, top=4pt, bottom=4pt, fonttitle=\bfseries, #1}

\newcommand{\E}{\mathbb{E}}
\newcommand{\Var}{\operatorname{var}}
\newcommand{\Cov}{\operatorname{cov}}
\renewcommand{\P}{\mathbb{P}}
\newcommand{\given}{\mid}
\newcommand{\Bern}{\text{Bernoulli}}
\newcommand{\Unif}{\text{Uniform}}
\newcommand{\Norm}{\text{Normal}}
\newcommand{\py}[1]{\texttt{#1}}


\setlength{\parindent}{0pt}
\emergencystretch=3em
\setlength{\parskip}{5pt}
\setlength{\headheight}{14pt}
\pagestyle{fancy}
\fancyhf{}
\lhead{Math 50 -- Midterm 1 study notes}
\rhead{Units 1--2}
\cfoot{\thepage}

\title{\textbf{Math 50 Midterm 1 Study Notes}\\[2pt]\large Unit 1 (probability models) and Unit 2 (expectation, Normal, regression)}
\author{}
\date{}

\begin{document}
\maketitle
\thispagestyle{fancy}

\begin{tipbox}[title={How to use these notes}]
Sections 1--2 are the concepts, in the same order as the ``What you need to know'' list on the exam page.
The last page is a one-page landscape formula sheet. The 21 practice problems and their solutions are in the separate file \texttt{practice\_solutions.pdf}. Everything here uses only the notation and tools from the course notes
(in particular, correlation is $\rho$).
\end{tipbox}

%==========================================================
\section*{0. Must-memorize checklist}
%==========================================================
\begin{center}
\renewcommand{\arraystretch}{1.35}
\begin{tabular}{@{}llll@{}}
\toprule
Distribution & Sample space / density & Mean & Variance\\
\midrule
$Y\sim\Bern(q)$ & $P(1)=q,\ P(0)=1-q$ & $q$ & $q(1-q)$\\
$Y\sim\Unif(a,b)$ & $f(y)=\dfrac{1}{b-a}$ on $[a,b]$ & $\dfrac{a+b}{2}$ & $\dfrac{(b-a)^2}{12}$\,$^*$\\
$Y\sim\Norm(\mu,\sigma^2)$ & $\dfrac{1}{\sqrt{2\pi\sigma^2}}e^{-(y-\mu)^2/2\sigma^2}$ & $\mu$ & $\sigma^2$\\
\bottomrule
\end{tabular}
\end{center}
{\small $^*$ The variance of a Uniform is \emph{not} derived in the notes (that would need an integral), but the exam
page says to memorize Uniform, and some practice problems need it. Bell curve: $68\%$ within $1\sigma$, $95\%$ within $2\sigma$, $99.7\%$ within $3\sigma$.}

Also know cold: \textbf{(i)} $\P(x,y)=\P(y\given x)\P(x)$ and Bayes; \textbf{(ii)} independence $\iff \P(x,y)=\P(x)\P(y)$;
\textbf{(iii)} $\E[Y]=\E[\E[Y\given X]]$; \textbf{(iv)} $\Var(Y)=\E[\Var(Y\given X)]+\Var(\E[Y\given X])$;
\textbf{(v)} $aX+b$ and sums of independent Normals; \textbf{(vi)} $\beta_1=\E[Y\given X=1]-\E[Y\given X=0]$;
\textbf{(vii)} $\Cov(X,Y)=\beta_1\sigma_X^2$.

%==========================================================
\section{Unit 1: Probability models}
%==========================================================

\subsection{Sample spaces, outcomes, events}
A random variable is a quantity we cannot predict before observing it (capital letters $X,Y,Z$). Its \textbf{sample space}
$S$ (or $S_X$) is the set of all values it can take. Elements of $S$ are \textbf{outcomes}; subsets of $S$ are \textbf{events}.
For an event $U$: $\P(U)=\sum_{x\in U}\P(x)$.

\begin{keybox}[title={Axioms}]
Nonnegativity $\P(U)\ge 0$; normalization $\P(S)=1$; additivity for disjoint events $\P(\cup U_i)=\sum\P(U_i)$.
\end{keybox}

\textbf{Examples.} Coin: $S=\{\text{heads},\text{tails}\}$ (coded $\{1,0\}$). Die: $S=\{1,\dots,6\}$. Two independent coin-type
variables $(X_1,X_2)$: $S=\{(0,0),(0,1),(1,0),(1,1)\}$. Compound variables (Drill 4): list every reachable value; a value can be reached
in more than one way (e.g.\ ``die $+1$ if heads, die $+4$ if tails'' gives $5$ from $(H,4)$ and $(T,1)$).

\subsection{The $\sim$ notation and the Bernoulli distribution}
``$Y\sim\Bern(q)$'' means $Y\in\{0,1\}$ with $\P(Y=1)=q$, $\P(Y=0)=1-q$. The probability of any event $E$ is the
probability of the Bernoulli variable $\mathbf 1_E$. General pattern: \emph{Variable} $\sim$ \emph{Distribution(parameters)}.

\subsection{Joint, marginal, conditional; independence}
\begin{itemize}[leftmargin=*,itemsep=2pt]
  \item \textbf{Marginal} (sum out the other variable): $\P(x)=\sum_y\P(x,y)$. On a table, add across a row or down a column.
  \item \textbf{Conditional}: $\P(x\given y)=\dfrac{\P(x,y)}{\P(y)}$ (restrict to the $y$ row/column, renormalize).
  \item \textbf{Product form}: $\P(x,y)=\P(x\given y)\P(y)$. \textbf{Bayes}: $\P(x\given y)=\dfrac{\P(y\given x)\P(x)}{\P(y)}$.
  \item Conditioning is \emph{not} symmetric: $\P(A\given B)\neq\P(B\given A)$ in general.
  \item \textbf{Independent} $\iff\P(x,y)=\P(x)\P(y)$ for all $x,y$ $\iff\P(y\given x)=\P(y)$ $\iff\P(x\given y)=\P(x)$.
\end{itemize}

\begin{tipbox}[title={Exam move: ``are $X$ and $Y$ independent? Justify with a single calculation''}]
To show \emph{not} independent you only need \emph{one} cell where $\P(x,y)\ne\P(x)\P(y)$ (or one conditional that differs from its marginal).
To show independent you must check every cell. Be careful when the numbers are close: compute exactly, do not eyeball.
\end{tipbox}

\textbf{Worked example (from the notes).} $\P(Y_A,Y_B)$: $(0,0)\!:\tfrac12$, $(0,1)\!:\tfrac18$, $(1,0)\!:\tfrac18$, $(1,1)\!:\tfrac14$.
Marginals: $\P(Y_A=0)=\tfrac58$, $\P(Y_B=0)=\tfrac58$. Then $\P(Y_A=0)\P(Y_B=0)=\tfrac{25}{64}\neq\tfrac12$, so not independent.
Conditional: $\P(Y_A=1\given Y_B=0)=\dfrac{1/8}{5/8}=\tfrac15$.

\subsection{From a $\sim$ model to the joint table, and back}
\textbf{Model $\to$ table.} A model like
\[
Y\sim\Bern(1/2),\qquad X\given Y\sim\Bern(Y/4+1/4)
\]
means: first draw $Y$; then, given $Y$, draw $X$ with success probability $Y/4+1/4$. So
$\P(X{=}1\given Y{=}0)=\tfrac14$, $\P(X{=}1\given Y{=}1)=\tfrac12$. Multiply along branches:
$\P(x,y)=\P(x\given y)\P(y)$.
\begin{center}
\begin{tabular}{c|cc|c}
$\P(x,y)$ & $Y=0$ & $Y=1$ &\\\hline
$X=0$ & $\tfrac12\cdot\tfrac34=\tfrac38$ & $\tfrac12\cdot\tfrac12=\tfrac14$ &\\
$X=1$ & $\tfrac12\cdot\tfrac14=\tfrac18$ & $\tfrac12\cdot\tfrac12=\tfrac14$ &\\\hline
 & $\tfrac12$ & $\tfrac12$ & $1$
\end{tabular}
\end{center}
\textbf{Table $\to$ model.} Read $Y$'s marginal off the column sums, then divide each column by its sum to get $X\given Y$, and recognize the
conditional as a named distribution whose parameter depends on $y$. Always check the table sums to $1$.

\subsection{Continuous distributions}
A continuous random variable has a \textbf{density} $f(y)\ge0$ with $\int f=1$, and probabilities are areas:
\[
\P(a<Y<b)=\int_a^b f(y)\,dy,\qquad \text{CDF: } F(y)=\P(Y\le y).
\]
\begin{itemize}[leftmargin=*,itemsep=2pt]
  \item $\P(Y=y)=0$ for every single point; only intervals have positive probability.
  \item \textbf{A density is not a probability}: $f(y)$ can exceed $1$ (e.g.\ $\Unif(0,0.5)$ has $f=2$); only $\int f\le 1$. Think of $f(y)\Delta y\approx\P(y\le Y\le y+\Delta y)$.
  \item $Y\sim\Unif(a,b)$: $f(y)=\frac1{b-a}$ on $[a,b]$, so $\P(y_1<Y<y_2)=\frac{y_2-y_1}{b-a}$ (proportional to length).
  \item Conditioning a Uniform on a sub-interval gives a Uniform on the sub-interval: $A\sim\Unif(20,90)$, $A\given(A<50)\sim\Unif(20,50)$.
  Equivalent recipe: $\P(Y<1\given Y<3)=\P(Y<1)/\P(Y<3)$ because $\{Y<1\}\subset\{Y<3\}$.
  \item No integrals will be required; Normal probabilities use the bell-curve rule (Unit 2).
\end{itemize}

\subsection{Sampling and simulation: translating code to probability statements}
A simulation approximates probability by frequency: $\P(X=x)\approx N(x)/n$. \textbf{Key idea:} a probability is the mean of a boolean array,
and conditioning means \emph{restricting the array first}, then averaging.

\begin{center}
\renewcommand{\arraystretch}{1.25}
\begin{tabular}{@{}p{0.5\linewidth}p{0.45\linewidth}@{}}
\toprule
Python & Probability meaning\\
\midrule
\py{X = np.random.choice(vals, p=p, size=N)} & $N$ iid draws with $\P(X=\text{vals}[i])=p[i]$\\
\py{np.mean(X == x)} & $\P(X=x)$\\
\py{np.mean((X == x) \& (Y == y))} & $\P(X=x,\,Y=y)$\\
\py{np.mean(X[Y == y] == x)} & $\P(X=x\given Y=y)$\\
\py{np.mean(X)} & $\E[X]$\\
\py{np.mean(X[Y == y])} & $\E[X\given Y=y]$\\
\py{np.var(X)} & $\Var(X)$\\
\py{np.random.normal(mu, sigma, n)} & draws from $\Norm(\mu,\sigma^2)$ (takes the \emph{standard deviation})\\
\py{np.random.uniform(a, b, n)} & draws from $\Unif(a,b)$\\
\py{np.cov(x, y)[0,1]} / \py{np.cov(x, y)[0,0]} & estimates $\Cov(X,Y)/\Var(X)=\beta_1$\\
\bottomrule
\end{tabular}
\end{center}

\textbf{Reading a simulation as a model.} An \py{if}/mask that draws $x$ with different parameters depending on $y$ is a conditional model $X\given Y$:
\begin{lstlisting}
y = np.random.choice([0, 1], p=[0.4, 0.6])        # Y ~ Bernoulli(0.6)
if y == 0:
    x = np.random.choice([0, 1], p=[0.8, 0.2])    # X | Y=0 ~ Bernoulli(0.2)
else:
    x = np.random.choice([0, 1], p=[0.3, 0.7])    # X | Y=1 ~ Bernoulli(0.7)
\end{lstlisting}
Model: $Y\sim\Bern(0.6)$, $X\given Y\sim\Bern(0.2+0.5Y)$. Joint: $\P(X{=}1,Y{=}0)=0.4(0.2)=0.08$, $\P(X{=}1,Y{=}1)=0.6(0.7)=0.42$,
so $\P(X{=}1)=0.50$. Independent? No, since $\P(X{=}1\given Y{=}0)=0.2\ne0.7=\P(X{=}1\given Y{=}1)$.

\textbf{Vectorized mask pattern (used on the exam page's practice problems).}
\begin{lstlisting}
x = np.random.choice([0, 1], size=n, p=[0.3, 0.7])
y = np.empty(n)
y[x == 0] = np.random.uniform(0, 2, len(y[x == 0]))   # Y | X=0 ~ Uniform(0,2)
y[x == 1] = np.random.uniform(1, 4, len(y[x == 1]))   # Y | X=1 ~ Uniform(1,4)
\end{lstlisting}
Each line assigns a \emph{conditional} distribution to the rows selected by the mask. Also, \py{np.random.normal(3 + 4*x, 2, n)} means
$Y\given X\sim\Norm(3+4X,\,2^2)$ (mean is an array, so it depends on each row's $x$; the $2$ is a standard deviation, so the variance is $4$).

\textbf{Predict-the-output rules of thumb.}
\begin{enumerate}[leftmargin=*,itemsep=2pt]
  \item Identify the model (in $\sim$ notation) first. Then the printed number is the corresponding exact quantity, up to sampling error ($\approx 1/\sqrt n$).
  \item \py{np.mean(x)} $\to\E[X]$ computed by the tower property (weight conditional means by group probabilities).
  \item \py{np.mean(x[y==1])} $\to\E[X\given Y=1]$: just the parameter used in that branch.
  \item \py{np.mean(y[x>0])} $\to\P(Y=1\given X>0)$ when $y$ is $0/1$ (a mean of $0/1$ values is a probability).
  \item Two arrays generated in separate lines given the same $x$ are \emph{conditionally independent given $x$}, but generally not independent marginally.
\end{enumerate}

%==========================================================
\section{Unit 2: Expectation, the Normal, and regression}
%==========================================================

\subsection{Expectation, variance, CV}
\[
\E[Y]=\sum_y y\,\P(y),\qquad \E[g(Y)]=\sum_y g(y)\P(y),\qquad
\Var(Y)=\E[(Y-\E Y)^2]=\E[Y^2]-(\E Y)^2 .
\]
Standard deviation $\sigma=\sqrt{\Var(Y)}$ (same units as $Y$). \textbf{Coefficient of variation} $\mathrm{CV}=\sigma/\mu$ (unitless: compares spread
across different scales; blows up when the mean is near $0$; shifting all values by a constant changes $\mu$ but not $\sigma$).

\begin{keybox}[title={Properties of expectation}]
\begin{enumerate}[leftmargin=*,itemsep=1pt]
  \item Linearity: $\E[X+Y]=\E[X]+\E[Y]$ (always, no independence needed). So $\E[aX+b]=a\E[X]+b$.
  \item $\Var(aX+b)=a^2\Var(X)$ (shifts don't matter, scale gets squared).
  \item Independent $\Rightarrow\E[XY]=\E[X]\E[Y]$, and $\Var(X+Y)=\Var(X)+\Var(Y)$.
  \item \textbf{Tower property:} $\E[Y]=\E[\E[Y\given X]]=\sum_x \E[Y\given X=x]\,\P(X=x)$.
  \item \textbf{Law of total variance:} $\Var(Y)=\underbrace{\E[\Var(Y\given X)]}_{\text{within-group}}+\underbrace{\Var(\E[Y\given X])}_{\text{between-group}}$.
\end{enumerate}
\end{keybox}

\textbf{Conditional expectation} $\E[Y\given X=x]=\sum_y y\,\P(y\given x)$: from data, average $Y$ over only those rows with $X=x$.
$\E[Y\given X]$ is itself a random variable (a function of $X$). Careful: $\E[Y\given X=x]=\E[Y]$ for all $x$ does \emph{not} imply independence.

\textbf{Bernoulli/Binomial facts.} $Y\sim\Bern(q)$: $\E Y=q$, $\Var Y=q(1-q)$ (since $Y^2=Y$). If $Y=\sum_{i=1}^N X_i$ with $X_i$ iid $\Bern(q)$:
$\E Y=Nq$, $\Var Y=Nq(1-q)$, $\mathrm{CV}=\sqrt{(1-q)/(qN)}\to0$.

\textbf{Worked example (tower + total variance).} Morning ($X{=}0$, prob $0.6$, mean $78$), afternoon ($X{=}1$, prob $0.4$, mean $86$), both with variance $64$:
\begin{align*}
\E[Y]&=78(0.6)+86(0.4)=81.2,\qquad \E[\Var(Y\given X)]=64,\\
\Var(\E[Y\given X])&=0.6(78-81.2)^2+0.4(86-81.2)^2=15.36,
\end{align*}
so $\Var(Y)=64+15.36=79.36$. Total variance is bigger than either group's, because the group means also differ.

\subsection{The Normal distribution}
$X\sim\Norm(\mu,\sigma^2)$ is symmetric, bell-shaped, determined by mean $\mu$ (center) and variance $\sigma^2$ (spread). \textbf{Watch the convention:}
math notation $\Norm(\mu,\sigma^2)$ takes the variance; \py{np.random.normal(mu, sigma)} takes the standard deviation.

\textbf{Standardizing.} $Z=\dfrac{X-\mu}{\sigma}\sim\Norm(0,1)$: measures $X$ in standard-deviation units.

\textbf{Bell-curve rule.}
\begin{center}
\begin{tabular}{@{}lccc@{}}
\toprule
Interval & $\mu\pm\sigma$ & $\mu\pm2\sigma$ & $\mu\pm3\sigma$\\
Probability & $\approx68\%$ & $\approx95\%$ & $\approx99.7\%$\\
\bottomrule
\end{tabular}
\end{center}
By symmetry, each tail beyond $1\sigma$ is $\approx16\%$ ($(100-68)/2$), beyond $2\sigma$ is $\approx2.5\%$, beyond $3\sigma$ is $\approx0.15\%$.
Also $\P(Z<1)\approx0.84$. \textbf{Recipe:} standardize the endpoints, then count how many $\sigma$'s they are from $\mu$.

\begin{keybox}[title={Linear transformations and sums}]
If $X\sim\Norm(\mu_x,\sigma_x^2)$ then $aX+b\sim\Norm(a\mu_x+b,\ a^2\sigma_x^2)$.\\
If $X_1\sim\Norm(\mu_1,\sigma_1^2)$, $X_2\sim\Norm(\mu_2,\sigma_2^2)$ \textbf{independent}:
\[
aX_1+bX_2+d\sim\Norm\big(a\mu_1+b\mu_2+d,\ a^2\sigma_1^2+b^2\sigma_2^2\big).
\]
Variances add (with squared coefficients) even for a difference: $\Var(X_1-X_2)=\sigma_1^2+\sigma_2^2$.
Average of $N$ independent copies of $\Norm(\mu,\sigma^2)$ is $\Norm(\mu,\sigma^2/N)$: same center, smaller spread.
\end{keybox}

\textbf{Example.} $Y\sim\Norm(5,4)$, so $\sigma=2$. $\P(Y>7)$: $7=\mu+\sigma$, so $\approx16\%$. $\P(Y>3\given Y<7)=\dfrac{\P(3<Y<7)}{\P(Y<7)}\approx\dfrac{0.68}{0.84}\approx0.81$.

\subsection{The single-predictor regression model}
\[
Y\given X\sim\Norm(\beta_0+\beta_1X,\ \sigma^2)\qquad\Longleftrightarrow\qquad Y=\beta_0+\beta_1X+\epsilon,\quad \epsilon\sim\Norm(0,\sigma^2).
\]
\begin{itemize}[leftmargin=*,itemsep=2pt]
  \item $\beta_0$: intercept, $\E[Y\given X=0]$ (the line's height at $x=0$).
  \item $\beta_1$: slope, the change in average $Y$ per one-unit increase in $X$.
  \item $\sigma^2$ ($=\sigma_\epsilon^2$): noise variance, the vertical spread of $Y$ around the line (same at every $x$).
\end{itemize}

\textbf{Assumptions.}
\begin{enumerate}[leftmargin=*,itemsep=1pt]
  \item \emph{Linearity}: $\E[Y\given X]=\beta_0+\beta_1X$ is a straight line.
  \item \emph{Homoskedasticity}: $\Var(Y\given X)=\sigma^2$ does not depend on $X$.
  \item \emph{Conditional normality}: $Y\given X$ (equivalently $\epsilon$) is Normal.
  \item \emph{Exogeneity}: $\E[\epsilon\given X]=0$, i.e.\ nothing left in the noise is related to $X$. If it fails, $X$ is \emph{endogenous}.
\end{enumerate}
\textbf{Spotting failures on a scatter plot:} curved cloud $\to$ linearity fails; spread that fans out or shrinks with $x$ $\to$ homoskedasticity fails;
strong skew or outliers in the vertical direction $\to$ normality fails. \textbf{Exogeneity cannot be checked from a scatter plot of $(x,y)$ alone}:
it is a statement about unobserved factors inside $\epsilon$, and the plot only shows how $Y$ relates to $X$ (the effect of $X$ and the effect of any confounder look identical).
Exogeneity is guaranteed by \emph{randomized assignment} of $X$ (coin flip), and threatened by self-selection.

\subsection{Binary predictor: $\beta_1$ is a difference of means}
If $X\in\{0,1\}$ (control/treatment), then $\E[Y\given X=0]=\beta_0$ and $\E[Y\given X=1]=\beta_0+\beta_1$, so
\[
\boxed{\ \beta_1=\E[Y\given X=1]-\E[Y\given X=0].\ }
\]
Estimates: $\hat\beta_0=\bar Y_{X=0}$, $\hat\beta_1=\bar Y_{X=1}-\bar Y_{X=0}$ (in code, \py{y[x==1].mean() - y[x==0].mean()}).
Relabeling ($X'=1-X$) swaps the groups: $\beta_0'=\beta_0+\beta_1$, $\beta_1'=-\beta_1$.
With a fair coin, $\E[Y]=\beta_0+\beta_1/2$ (tower property).

\subsection{Covariance and the slope}
$\Cov(X,Y)=\E[XY]-\E[X]\E[Y]$, and $\Cov(X,X)=\Var(X)$.
\begin{keybox}[title={Key derivation: $\Cov(X,Y)=\beta_1\sigma_X^2$}]
Using $\E[XY]=\E[X\,\E[Y\given X]]=\E[X(\beta_0+\beta_1X)]=\beta_0\E[X]+\beta_1\E[X^2]$ and $\E[X^2]=\sigma_X^2+\mu_X^2$:
\[
\Cov(X,Y)=\beta_0\mu_X+\beta_1(\sigma_X^2+\mu_X^2)-\mu_X(\beta_0+\beta_1\mu_X)=\beta_1\sigma_X^2.
\]
Hence $\beta_1=\Cov(X,Y)/\Var(X)$, valid for any $X$ (not just binary), and
$\Var(Y)=\beta_1^2\sigma_X^2+\sigma_\epsilon^2$ (total variance: within $=\sigma_\epsilon^2$, between $=\beta_1^2\sigma_X^2$).
\end{keybox}
Marginal of $Y$ when $X$ is Normal: $Y\sim\Norm(\beta_0+\beta_1\mu_X,\ \beta_1^2\sigma_X^2+\sigma_\epsilon^2)$ (sum of independent Normals).
When $X$ is \emph{not} Normal (e.g.\ Bernoulli), $Y$ still has the same mean and variance formulas, but $Y$ is a mixture and is \emph{not} Normal.

\subsection{Correlation, $R^2$, regression to the mean}
\[
\rho=\frac{\beta_1\sigma_X}{\sigma_Y}=\frac{\Cov(X,Y)}{\sigma_X\sigma_Y},\qquad
\E[Z_Y\given Z_X]=\rho Z_X,\qquad
\rho^2=1-\frac{\sigma_\epsilon^2}{\sigma_Y^2}=R^2,\qquad -1\le\rho\le1 .
\]
$\rho$ is the regression slope after standardizing both variables (unitless); $\rho^2$ is the fraction of $\Var(Y)$ explained by $X$
($\rho^2=1$: $Y$ is a deterministic line in $X$; $\rho=0$: $X$ tells you nothing about the mean of $Y$).
\textbf{Regression to the mean:} if $|\rho|<1$, an unusually extreme $X$ (say $z$ SDs from its mean) predicts a \emph{less} extreme $Y$ ($\rho z$ SDs);
this is not a causal effect, just luck not repeating.
Rescaling $Y$ (e.g.\ units) changes $\beta_1$ and $\Cov$ but not $\rho$ or $R^2$.

\subsection{Exogeneity and randomized experiments (conceptual)}
Random assignment (a fair coin) makes $X$ independent of everything else, so $\E[\epsilon\given X]=0$ and $\hat\beta_1$ can be read as the effect of the treatment.
If subjects self-select into groups, unmeasured traits (motivation, health) that affect $Y$ also differ between groups: $\E[\epsilon\given X]\neq0$, $X$ is endogenous,
and $\hat\beta_1$ mixes the treatment effect with those differences.

%==========================================================
% Cheat sheet: landscape, 3 columns, one page
%==========================================================
\clearpage
\newgeometry{margin=0.4in}
\begin{landscape}
\pagestyle{empty}
\newcommand{\csh}[1]{\par\vspace{5pt}{\color{accent}\bfseries #1}\vspace{-1pt}\par{\color{accent}\hrule height 0.6pt}\vspace{2pt}\par}
\newtcolorbox{cbox}{colback=tipbg, colframe=orange!70!black, boxrule=0.5pt, arc=1.5pt,
  left=3pt, right=3pt, top=2pt, bottom=2pt, before skip=3pt, after skip=3pt}
\begingroup
\footnotesize
\setlength{\parskip}{2.5pt}
\setlength{\tabcolsep}{4pt}
{\Large\bfseries\color{accent} Math 50 Midterm 1 \textperiodcentered\ Cheat sheet}\hfill{\small Units 1--2}\par
{\color{accent}\hrule height 1pt}\vspace{2pt}
\begin{multicols}{3}
\raggedright

\csh{Named distributions}
\renewcommand{\arraystretch}{1.25}
\begin{tabular}{@{}lccc@{}}
 & pmf / density & mean & var\\\hline
$\Bern(q)$ & $\P(0)=1-q,\ \P(1)=q$ & $q$ & $q(1-q)$\\
$\Unif(a,b)$ & $\tfrac1{b-a}$ on $[a,b]$ & $\tfrac{a+b}2$ & $\tfrac{(b-a)^2}{12}$\\
$\Norm(\mu,\sigma^2)$ & (below) & $\mu$ & $\sigma^2$\\
\end{tabular}

Normal density $\dfrac1{\sqrt{2\pi\sigma^2}}e^{-(x-\mu)^2/(2\sigma^2)}$; second parameter is the \emph{variance}.

\begin{tabular}{@{}lccc@{}}
Bell curve & $\mu\pm\sigma$ & $\mu\pm2\sigma$ & $\mu\pm3\sigma$\\\hline
inside & 68\% & 95\% & 99.7\%\\
one tail beyond & 16\% & 2.5\% & 0.15\%\\
\end{tabular}

$\P(Z<1)\approx84\%$, $\P(Z<0)=50\%$.

\csh{Probability}
\textbf{Events:} $\P(U)=\sum_{x\in U}\P(x)$, $\P(S)=1$, $\P(\text{not }U)=1-\P(U)$; ``at least one'' $=1-\P(\text{none})$.\\
\textbf{Marginal:} $\P(x)=\sum_y\P(x,y)$; through a conditional: $\P(y)=\sum_x\P(y\given x)\P(x)$.\\
\textbf{Conditional} ($\P(y)>0$): $\P(x\given y)=\dfrac{\P(x,y)}{\P(y)}$, so $\P(x,y)=\P(x\given y)\P(y)$.\\
\textbf{Bayes:} $\P(x\given y)=\dfrac{\P(y\given x)\P(x)}{\P(y)}$; $\P(A\given B)\ne\P(B\given A)$.\\
\textbf{$X,Y$ independent} $\iff\P(x,y)=\P(x)\P(y)$ for \emph{all} $x,y$ $\iff\P(y\given x)=\P(y)\iff\P(x\given y)=\P(x)$.\\
\textbf{Continuous:} density $f\ge0$, $\int f=1$; CDF $F(y)=\P(Y\le y)$; $\P(a<Y<b)=\int_a^bf=F(b)-F(a)$; $\P(Y{=}y)=0$; $f$ can exceed 1. Sub-interval: $\P(Y<1\given Y<3)=\P(Y<1)/\P(Y<3)$.

\csh{$\sim$ notation $\leftrightarrow$ joint table}
$Y\sim\Bern(q)$, $X\given Y\sim\Bern(g(Y))$:\\
$\P(X{=}1,Y{=}y)=g(y)\P(y)$, \ $\P(X{=}0,Y{=}y)=(1-g(y))\P(y)$, \ $\P(X{=}1)=\sum_yg(y)\P(y)$.\\
Table $\to$ model: sum over $x$ to get $\P(y)$; divide each $y$-slice by $\P(y)$ to get $X\given Y{=}y$. Check the table sums to 1.

\begin{cbox}
\textbf{Exam moves.}\\
$\bullet$ Independence: one failing cell $\Rightarrow$ dependent.\\
$\bullet$ Simulation: write the model in $\sim$ form first; a printed mean $\approx$ the exact $\E$ (tower).\\
$\bullet$ Mean of a 0/1 array is a probability.\\
$\bullet$ Sample average of rows meeting a condition $=$ conditional expectation.\\
$\bullet$ Densities can exceed 1; mixtures of Normals are not Normal.
\end{cbox}

\columnbreak

\csh{Expectation and variance}
$\E[g(Y)]=\sum_y g(y)\P(y)$; \ $\E[Y\given x]=\sum_y y\,\P(y\given x)$ (data: average over rows with $X{=}x$).\\
$\Var Y=\E[Y^2]-(\E Y)^2$, \ $\sigma=\sqrt{\Var Y}$; \ $\Var(Y\given x)=\E[Y^2\given x]-\E[Y\given x]^2$.\\
$\E[Y^2]\ne(\E Y)^2$: $\E[Y^2]=\Var Y+(\E Y)^2\ge(\E Y)^2$, equal only if $\Var Y=0$. Compute $\E[Y^2]=\sum_y y^2\P(y)$ (square first, then average). Bernoulli: $Y^2=Y$, so $\E[Y^2]=q$.\\
$\mathrm{CV}=\sigma/\mu$: unitless; shifting values changes $\mu$, not $\sigma$; blows up near $\mu=0$.\\
$\E[aX+b]=a\E X+b$, \quad $\Var(aX+b)=a^2\Var X$.\\
$\E[X+Y]=\E X+\E Y$ always. If independent: $\E[XY]=\E X\,\E Y$ and $\Var(X+Y)=\Var X+\Var Y$.\\
Binomial ($N$ iid $\Bern(q)$): $\E=Nq$, $\Var=Nq(1-q)$, $\mathrm{CV}=\sqrt{\tfrac{1-q}{qN}}$.

\begin{cbox}
\textbf{Tower:} $\E Y=\sum_x\E[Y\given x]\,\P(x)$.\\
\textbf{Total variance:}\\
$\Var Y=\E[\Var(Y\given X)]+\Var(\E[Y\given X])$\\
{\scriptsize (within-group + between-group)}
\end{cbox}

\csh{Normal}
\textbf{Standardize:} $Z=\dfrac{X-\mu}{\sigma}\sim\Norm(0,1)$. Recipe: standardize the endpoints, count $\sigma$'s.\\
$aX+b\sim\Norm(a\mu+b,\ a^2\sigma^2)$.\\
Independent sum:\\
$aX_1+bX_2+d\sim\Norm\big(a\mu_1+b\mu_2+d,\ a^2\sigma_1^2+b^2\sigma_2^2\big)$; a difference still \emph{adds} variances.\\
Mean of $N$ iid: $\Norm(\mu,\sigma^2/N)$.\\
$\P(X>\mu)=\tfrac12$ (symmetry).\\
A mixture of Normals with different means (e.g.\ $Y\given X$ with $X$ Bernoulli) is \emph{not} Normal.

\csh{Regression model}
$Y\given X\sim\Norm(\beta_0+\beta_1X,\sigma^2)\iff Y=\beta_0+\beta_1X+\epsilon$, $\epsilon\sim\Norm(0,\sigma^2)$.\\
$\beta_0=\E[Y\given X{=}0]$; \ $\beta_1$ = change in mean $Y$ per unit $X$; \ $\sigma^2=\Var(Y\given X)$ = vertical spread, same at every $x$.

\textbf{Assumptions:} \textbf{(1)} linear mean; \textbf{(2)} constant variance; \textbf{(3)} Normal errors; \textbf{(4)} exogeneity $\E[\epsilon\given X]=0$.\\
Fails on a plot: curve $\to$ (1); fanning spread $\to$ (2); skew/outliers $\to$ (3). \textbf{(4) cannot be checked from a scatter plot.} Randomization $\Rightarrow$ exogenous; self-selection $\Rightarrow$ maybe endogenous.

\columnbreak

\csh{Binary predictor $X\in\{0,1\}$}
$\beta_1=\E[Y\given X{=}1]-\E[Y\given X{=}0]$\\
$\hat\beta_0=\bar Y_{X=0}$, \quad $\hat\beta_1=\bar Y_{X=1}-\bar Y_{X=0}$\\
$\E Y=\beta_0+\beta_1\P(X{=}1)$\\
$X\sim\Bern(p)$: $\sigma_X^2=p(1-p)$ ($\tfrac14$ if fair), so $\Cov=\beta_1p(1-p)$.\\
Relabel $X'=1-X$: $\beta_0'=\beta_0+\beta_1$, $\beta_1'=-\beta_1$.

\csh{Covariance, correlation, $R^2$}
$\Cov(X,Y)=\E[XY]-\E X\,\E Y$, \ $\Cov(X,X)=\Var X$
\begin{cbox}
$\Cov(X,Y)=\beta_1\sigma_X^2\ \Rightarrow\ \beta_1=\dfrac{\Cov(X,Y)}{\Var X}$\\
{\scriptsize any $X$ with $\Var X<\infty$}
\end{cbox}
$\E Y=\beta_0+\beta_1\E X$; \ $\Var Y=\beta_1^2\sigma_X^2+\sigma_\epsilon^2$ (between + within).\\
$X$ Normal: $Y\sim\Norm(\beta_0+\beta_1\mu_X,\ \beta_1^2\sigma_X^2+\sigma_\epsilon^2)$. $X$ not Normal: same mean and variance, but $Y$ is not Normal.\\
$\rho=\dfrac{\beta_1\sigma_X}{\sigma_Y}=\dfrac{\Cov(X,Y)}{\sigma_X\sigma_Y}\in[-1,1]$, sign of $\rho$ = sign of $\beta_1$.\\
$R^2=\rho^2=1-\dfrac{\sigma_\epsilon^2}{\sigma_Y^2}$; \ $\E[Z_Y\given Z_X]=\rho Z_X$ (regression to the mean).\\
Changing units of $Y$ changes $\beta_1$, $\Cov$; not $\rho$, $R^2$.\\
Independent $\Rightarrow\Cov=0$, but $\Cov=0\not\Rightarrow$ independent.

\csh{Code $\to$ math}
\renewcommand{\arraystretch}{1.2}
\begin{tabular}{@{}p{0.62\linewidth}p{0.33\linewidth}@{}}
\py{np.random.choice(vals, p=p, size=n)} & $n$ iid draws from pmf $p$\\
\py{np.random.uniform(a, b, n)} & $\Unif(a,b)$\\
\py{np.random.normal(mu, sigma, n)} & $\Norm(\mu,\sigma^2)$, \emph{SD} input\\
\py{np.random.normal(3+4*x, 2, n)} & $\Norm(3+4X,\,2^2)$\\
\py{y[x==0] = ...} & $Y\given X{=}0$\\
\py{np.mean(X==x)} & $\P(X{=}x)$\\
\py{np.mean((X==x) \& (Y==y))} & $\P(X{=}x,Y{=}y)$\\
\py{np.mean(X[Y==y]==x)} & $\P(x\given y)$\\
\py{np.mean(X)} / \py{np.var(X)} & $\E X$ / $\Var X$\\
\py{np.mean(X[Y==y])} & $\E[X\given y]$\\
\py{y[x==1].mean() - y[x==0].mean()} & $\approx\beta_1$ (binary $x$)\\
\py{np.cov(x,y)[0,1] / np.cov(x,y)[0,0]} & $\approx\beta_1$\\
\end{tabular}

\end{multicols}
\endgroup
\end{landscape}
\restoregeometry

\end{document}
